% Mordell--Weil heights explode: on y^2=x^3-2 with generator P=(3,5), the numerator of
% x(nP) has 1,3,6,10,15 digits for n=1..5 -- the canonical height grows like n^2 h(P).
\begin{tikzpicture}[>=Latex]
  \draw[->, black!55, line width=0.6pt] (0,0)--(6.4,0) node[right, font=\small] {$n$};
  \draw[->, black!55, line width=0.6pt] (0,0)--(0,8.4)
    node[above, font=\small, align=center] {height};
  \foreach \n/\v in {1/1, 2/3, 3/6, 4/10, 5/15}{
    \fill[formteal!55!black] (\n-0.30,0) rectangle (\n+0.30,{\v*0.5});
    \node[font=\small, below] at (\n,0) {$\n$};
    \node[font=\footnotesize, formteal!50!black, above=1pt] at (\n,{\v*0.5}) {$\v$};}
  \draw[formrust!78!black, line width=1.3pt, smooth, samples=40, domain=0.35:5.25]
    plot (\x,{0.30*\x*\x});
  \node[formrust!72!black, font=\small, anchor=west] at (1.6,7.4) {$\hat h(nP)=n^2\,\hat h(P)$};
\end{tikzpicture}
